\documentclass[11pt,spanish]{article}
%\usepackage[latin1]{inputenc}
\usepackage[utf8]{inputenc}
\usepackage[spanish,es-nodecimaldot]{babel}
\usepackage{graphicx}
\usepackage{tikz}
\usepackage{amsmath,amssymb}

\oddsidemargin=0.5cm
\textwidth=16cm
\textheight=48\baselineskip
\topmargin=-1.5cm
\parskip 2truemm
\parindent 0pt
\newcounter{problema}
\typeout{****************************************}
\typeout{****** Configuracion de Guia 11pt ******}
\typeout{****************************************}
%%%%%%%%%%%%%%%%%%%%%%%%%%%%%%%%%%%%%%%%%%%%%%%%%%%%%%%%%%%%%%%%%%%%%%%

\begin{document}

\renewcommand{\labelenumi}{\em \alph{enumi})}
\renewcommand{\labelenumii}{\em \arabic{enumii})}
\newcommand\al{\alpha}
\newcommand\n{\nabla}
\newcommand\F{\vec F}
\newcommand\G{\vec G}
\newcommand\noi{\noindent}
\newcommand\re{\mbox{Re}}
\newcommand\im{\mbox{Im}}
%\newcommand\arg{\mbox{arg}}
\newcommand\Arg{\mbox{Arg}}
\newcommand{\prob}{\stepcounter{problema} \noindent{\bf Problema
\arabic{problema}: }}
\newcommand\pin[2]{\big\langle #1,#2 \big\rangle}

%\input{/home/ortiz/lib/definiciones.tex}
\begin{center}
{\large\sc Facultad de Matem\'atica, Astronom\'ia, F\'isica y Computaci\'on, U.N.C.}

\vspace{2mm}
{\large M\'etodos Matem\'aticos de la F\'isica II}

\vspace{2mm}
{\large Gu\'ia N$^{\sf o}$ 5 (2021\color[rgb]{0,0,0})}
\end{center}

\prob{} Considere la ecuaci\'on de calor $\kappa \Delta u= \cfrac{\partial u}{\partial t}$ dentro de una caja delimitada por $0<x<L$, $0 < y <H$ y $0 < z <W$.  En $t=0$ la distribuci\'on de tempratura es $u=2y \sen(x)$.
Las condiciones de borde son:

\begin{tabular}{ccc}
$u(0,y,z,t)=0$ \hspace{1cm} & $\cfrac{\partial u}{\partial y} (x,0,z,t)=0$ \hspace{1cm} & $\cfrac{\partial u}{\partial z} (x,y,0,t)=0$ \\
$u(L,y,z,t)=0$ \hspace{1cm} & $\cfrac{\partial u}{\partial y} (x,H,z,t)=0$ \hspace{1cm} & $u(x,y,W,t)=0$ \\

\end{tabular}

\prob{} Resuelva la ecuación de calor $\kappa \Delta u= \cfrac{\partial u}{\partial t}$ dentro de un cilindro de radio $a$ y altura $H$ con condiciones de borde $u(r,\theta,0,t)=0$, $u(r,\theta,H,t)=0$ y $u(a,\theta,z,t)=0$. Inicialmente el cilindro tiene una distribución de temperatua dada por $u(r,\theta,z,0)=rz$.

\prob{} Encuentre la distribuci\'on de temperatura, $u(r,\theta,\phi,t)$,   dentro de una esfera de radio $a$ si $u(a,\theta,\phi,t)=0$ y $u(r,\theta,\phi,0)=\sen(\theta)$, donde $\theta \in [0, \pi]$.

\prob{} Considere una placa metálica semicircular de radio $a$, es decir $0 < \theta < \pi$ y $0 < r < a$. Resuelva la ecuación de calor $\kappa \Delta u= \cfrac{\partial u}{\partial t}$ en el semicírculo, con las condiciones de contorno de Dirichlet:
$u(r,0,t) = 0$, $u(r,\pi,t) = 0$, $u(a,\theta,t) = 0$, y la condición inicial $u(r,\theta,0) = r \sen^2(\theta)$.

\prob{} Resuelva la ecuación de calor $\kappa \Delta u= \cfrac{\partial u}{\partial t}$ dentro de un cuarto de círculo, con las condiciones de contorno: $u(r,0,t) = 0$, $u(r,\pi/2,t) = 0$, $\cfrac{\partial u}{\partial r}\,(R,\theta,t) = 0$, y la condición inicial $u(r,\theta,0) = r^2 \cos(3\theta)$.

\prob{} Considere la ecuación de calor unidimensional con un término fuente
\[\cfrac{\partial u}{\partial t} = \cfrac{\partial^2 u}{\partial x^2} + Q(x),\]
siendo $0 < x < L$ y $t > 0$, con las condiciones de contorno: $u(0,t) = 0$, $u(L,t) = 0$ y la condición inicial $u(x,0) = f(x)$. Resuelva dicha ecuación para el caso en que $L = 2$, $f(x) = 2x-x^2$, $Q(x) = 1 - |x-1|$.

\prob{} Resuelva la ecuación
\[\cfrac{\partial u}{\partial t} = \cfrac{\partial^2 u}{\partial x^2} + \alpha\,u,\]
siendo $0 < x < 2$ y $t > 0$, con las condiciones de contorno: $\cfrac{\partial u}{\partial x}\,(0,t) = 0$, $\cfrac{\partial u}{\partial x}\,(2,t) = 0$, y la condición inicial $u(x,0) = 4x-x^3$.
\emph{Ayuda}: Realice el cambio de variable $u = A(x,t)\,v$ y determine la función $A(x,t)$ que lleve a resolver una ecuación de calor para $v(x,t)$ sin ningún término fuente (como $\alpha\,u$).


\end{document}


